% GENERATED FILE. Edit canonical lessons, then run npm run generate:books.
% canonical-source-hash: fnv1a64:921ac13a9f23f89b
% canonical-lessons: MW-C179-sil, MW-C179-kutno, MW-C179-dalno, MW-C179-chhilno, MW-C179-phatakno

\chapter{Grain and Spice}
\label{ch:mw-grain-and-spice}
\hlchaptermodality{179}

\begin{chapteropening}
\textbf{By the end of this chapter:} \emph{I can say how grain is cleaned and ground and spice pounded.}

Complete the last lesson of chapter 179: I can say how grain is cleaned and ground and spice pounded.
\end{chapteropening}

\section[\texorpdfstring{\mw{सिल} (sil)}{सिल (sil)}]{\mw{सिल} (sil) — a grinding stone}

\label{lesson:MW-C179-sil}



\begin{quote}
Before the new one: say the Marwadi for to scatter, then the Marwadi for to hand over.
\end{quote}

\subsection*{You'll want to know: \mw{सिल}}
\textbf{\mw{सिल}} — \emph{sil} — \textquotedblleft{}a grinding stone\textquotedblright{}.

\begin{quote}
\mw{रसोई} \mw{मांय} \mw{एक} \mw{भारी} \mw{सिल} \mw{है।}
\end{quote}

\begin{quote}
\emph{rasoī mā̃y ek bhārī sil hai.}
\end{quote}

\begin{quote}
There is a heavy grinding stone in the kitchen.
\end{quote}

A feminine noun.

\begin{cousinweb}
\textbf{\mw{सिल}} (\emph{sil}) is inherited from Sanskrit \emph{śilā}, \textquotedblleft{}a rock\textquotedblright{}, also \textquotedblleft{}the lower millstone\textquotedblright{}; CDIAL lists Hindi \emph{sil}, \textquotedblleft{}a flat grinding stone\textquotedblright{} (CDIAL 12459). CDIAL gives no Marwari form in this entry.
\end{cousinweb}

\begin{grammarlens}[title={What you've built}]
One more word for this chapter.
\end{grammarlens}

\subsection*{Your turn}
\begin{itemize}
\raggedright
  \item \emph{Say it:} \emph{sil}
  \item \emph{Say it:} \emph{sil}, once more
  \item \emph{Say it:} \emph{rasoī mā̃y ek bhārī sil hai.} — There is a heavy grinding stone in the kitchen.
  \item \emph{Recall:} say the Marwadi for to scatter, then the Marwadi for to hand over, then say \emph{sil} again
\end{itemize}

\begin{tcolorbox}[breakable,colback=teal!4,colframe=teal!35!black,title={Before you move on}]
What does \mw{सिल} mean? (\textquotedblleft{}A grinding stone\textquotedblright{}.) Say it once more.
\end{tcolorbox}
\section[\texorpdfstring{\mw{कूटणो} (kūṭṇo)}{कूटणो (kūṭṇo)}]{\mw{कूटणो} (kūṭṇo) — to pound}

\label{lesson:MW-C179-kutno}



\begin{quote}
Before the new one: say the Marwadi for to hand over, then the Marwadi for a grinding stone.
\end{quote}

\subsection*{You'll want to know: \mw{कूटणो}}
\textbf{\mw{कूटणो}} — \emph{kūṭṇo} — \textquotedblleft{}to pound\textquotedblright{}.

\begin{quote}
\mw{दादी} \mw{सिल} \mw{माथै} \mw{मिरच} \mw{कूटै} \mw{है।}
\end{quote}

\begin{quote}
\emph{dādī sil māthai mirach kūṭai hai.}
\end{quote}

\begin{quote}
Grandmother pounds chillies on the grinding stone.
\end{quote}

Past (I …): \mw{म्हैं} \mw{मिरच} \mw{कूटी}, \emph{mhaĩ mirach kūṭī}.

\begin{cousinweb}
\textbf{\mw{कूटणो}} (\emph{kūṭṇo}) is inherited from Sanskrit \emph{kuṭṭayati}, \textquotedblleft{}crushes, pounds\textquotedblright{}; CDIAL lists Hindi \emph{kūṭnā} and Gujarati \emph{kuṭvũ} (CDIAL 3241). CDIAL gives no Marwari form in this entry.
\end{cousinweb}

\begin{grammarlens}[title={What you've built}]
One more word for this chapter.
\end{grammarlens}

\subsection*{Your turn}
\begin{itemize}
\raggedright
  \item \emph{Say it:} \emph{kūṭṇo}
  \item \emph{Say it:} \emph{kūṭṇo}, once more
  \item \emph{Say it:} \emph{dādī sil māthai mirach kūṭai hai.} — Grandmother pounds chillies on the grinding stone.
  \item \emph{Recall:} say the Marwadi for to hand over, then the Marwadi for a grinding stone, then say \emph{kūṭṇo} again
\end{itemize}

\begin{tcolorbox}[breakable,colback=teal!4,colframe=teal!35!black,title={Before you move on}]
What does \mw{कूटणो} mean? (\textquotedblleft{}To pound\textquotedblright{}.) Say it once more.
\end{tcolorbox}
\section[\texorpdfstring{\mw{दळणो} (daḷṇo)}{दळणो (daḷṇo)}]{\mw{दळणो} (daḷṇo) — to grind (grain)}

\label{lesson:MW-C179-dalno}



\begin{quote}
Before the new one: say the Marwadi for a grinding stone, then the Marwadi for to pound.
\end{quote}

\subsection*{You'll want to know: \mw{दळणो}}
\textbf{\mw{दळणो}} — \emph{daḷṇo} — \textquotedblleft{}to grind (grain)\textquotedblright{}.

\begin{quote}
\mw{मां} \mw{चक्की} \mw{माथै} \mw{बाजरो} \mw{दळै} \mw{है।}
\end{quote}

\begin{quote}
\emph{mā̃ chakkī māthai bājro daḷai hai.}
\end{quote}

\begin{quote}
Mother grinds millet on the hand mill.
\end{quote}

Past (I …): \mw{म्हैं} \mw{बाजरो} \mw{दळ्यो}, \emph{mhaĩ bājro daḷyo}.

\begin{cousinweb}
\textbf{\mw{दळणो}} (\emph{daḷṇo}) is inherited from Sanskrit \emph{dalati}, \textquotedblleft{}cracks, splits\textquotedblright{}; CDIAL lists Hindi \emph{dalnā}, \textquotedblleft{}to grind coarsely\textquotedblright{}, and Gujarati \emph{daḷvũ} (CDIAL 6216). CDIAL gives no Marwari form in this entry.
\end{cousinweb}

\begin{grammarlens}[title={What you've built}]
One more word for this chapter.
\end{grammarlens}

\subsection*{Your turn}
\begin{itemize}
\raggedright
  \item \emph{Say it:} \emph{daḷṇo}
  \item \emph{Say it:} \emph{daḷṇo}, once more
  \item \emph{Say it:} \emph{mā̃ chakkī māthai bājro daḷai hai.} — Mother grinds millet on the hand mill.
  \item \emph{Recall:} say the Marwadi for a grinding stone, then the Marwadi for to pound, then say \emph{daḷṇo} again
\end{itemize}

\begin{tcolorbox}[breakable,colback=teal!4,colframe=teal!35!black,title={Before you move on}]
What does \mw{दळणो} mean? (\textquotedblleft{}To grind (grain)\textquotedblright{}.) Say it once more.
\end{tcolorbox}
\section[\texorpdfstring{\mw{छीलणो} (chhīlṇo)}{छीलणो (chhīlṇo)}]{\mw{छीलणो} (chhīlṇo) — to peel}

\label{lesson:MW-C179-chhilno}



\begin{quote}
Before the new one: say the Marwadi for to pound, then the Marwadi for to grind.
\end{quote}

\subsection*{You'll want to know: \mw{छीलणो}}
\textbf{\mw{छीलणो}} — \emph{chhīlṇo} — \textquotedblleft{}to peel\textquotedblright{}.

\begin{quote}
\mw{मां} \mw{आलू} \mw{छीलै} \mw{है।}
\end{quote}

\begin{quote}
\emph{mā̃ ālū chhīlai hai.}
\end{quote}

\begin{quote}
Mother peels the potatoes.
\end{quote}

Past (I …): \mw{म्हैं} \mw{केळो} \mw{छील्यो}, \emph{mhaĩ keḷo chhīlyo}.

\begin{cousinweb}
\textbf{\mw{छीलणो}} (\emph{chhīlṇo}) comes from a reconstructed \emph{*chilla}, \textquotedblleft{}torn, cut\textquotedblright{}; CDIAL lists Hindi \emph{chīlnā} and Gujarati \emph{chilvũ}, \textquotedblleft{}to peel\textquotedblright{} (CDIAL 5051). CDIAL gives no Marwari form in this entry.
\end{cousinweb}

\begin{grammarlens}[title={What you've built}]
One more word for this chapter.
\end{grammarlens}

\subsection*{Your turn}
\begin{itemize}
\raggedright
  \item \emph{Say it:} \emph{chhīlṇo}
  \item \emph{Say it:} \emph{chhīlṇo}, once more
  \item \emph{Say it:} \emph{mā̃ ālū chhīlai hai.} — Mother peels the potatoes.
  \item \emph{Recall:} say the Marwadi for to pound, then the Marwadi for to grind, then say \emph{chhīlṇo} again
\end{itemize}

\begin{tcolorbox}[breakable,colback=teal!4,colframe=teal!35!black,title={Before you move on}]
What does \mw{छीलणो} mean? (\textquotedblleft{}To peel\textquotedblright{}.) Say it once more.
\end{tcolorbox}
\section[\texorpdfstring{\mw{फटकणो} (phaṭakṇo)}{फटकणो (phaṭakṇo)}]{\mw{फटकणो} (phaṭakṇo) — to winnow}

\label{lesson:MW-C179-phatakno}



\begin{quote}
Before the new one: say the Marwadi for to grind, then the Marwadi for to peel.
\end{quote}

\subsection*{You'll want to know: \mw{फटकणो}}
\textbf{\mw{फटकणो}} — \emph{phaṭakṇo} — \textquotedblleft{}to winnow\textquotedblright{}.

\begin{quote}
\mw{मां} \mw{बाजरो} \mw{फटकै} \mw{है।}
\end{quote}

\begin{quote}
\emph{mā̃ bājro phaṭkai hai.}
\end{quote}

\begin{quote}
Mother winnows the millet.
\end{quote}

Past (I …): \mw{म्हैं} \mw{बाजरो} \mw{फटक्यो}, \emph{mhaĩ bājro phaṭakyo}.

\begin{cousinweb}
\textbf{\mw{फटकणो}} (\emph{phaṭakṇo}) comes from a reconstructed \emph{*phaṭ}, \textquotedblleft{}a sudden movement\textquotedblright{}, a word that imitates a sharp sound; CDIAL lists Hindi \emph{phaṭaknā}, \textquotedblleft{}to winnow\textquotedblright{} (CDIAL 9038). CDIAL gives no Marwari form in this entry.
\end{cousinweb}

\begin{grammarlens}[title={What you've built}]
One more word for this chapter.
\end{grammarlens}

\subsection*{Your turn}
\begin{itemize}
\raggedright
  \item \emph{Say it:} \emph{phaṭakṇo}
  \item \emph{Say it:} \emph{phaṭakṇo}, once more
  \item \emph{Say it:} \emph{mā̃ bājro phaṭkai hai.} — Mother winnows the millet.
  \item \emph{Recall:} say the Marwadi for to grind, then the Marwadi for to peel, then say \emph{phaṭakṇo} again
\end{itemize}

\begin{tcolorbox}[breakable,colback=teal!4,colframe=teal!35!black,title={Before you move on}]
What does \mw{फटकणो} mean? (\textquotedblleft{}To winnow\textquotedblright{}.) Say it once more.
\end{tcolorbox}
